% ==================================================
% Telescoping Identity
% Discrete Calculus
% ==================================================

\section*{The Telescoping Identity}

\begin{tcolorbox}[colback=thmbox, colframe=thmborder, arc=2pt,
  left=6pt, right=6pt, top=4pt, bottom=4pt,
  title={\small\textbf{Theorem (Telescoping Identity)}},
  fonttitle=\small\bfseries]
For any function \(f : \mathbb{Z} \to \mathbb{R}\) and integers \(a \le b\),
\[
\sum_{k=a}^{b-1} (f(k+1) - f(k)) = f(b) - f(a).
\]
\end{tcolorbox}
\label{thm:telescoping}

\begin{remark*}[Fully quantified form]
\(\forall f : \mathbb{Z}\to\mathbb{R},\; \forall a,b \in \mathbb{Z},\; a \le b:\;
\displaystyle\sum_{k=a}^{b-1}(f(k+1)-f(k)) = f(b)-f(a)\).
\end{remark*}

\begin{proof}
Expanding the sum gives
\[
\begin{aligned}
(f(a+1)-f(a)) &+ (f(a+2)-f(a+1)) \\
              &+ \cdots + (f(b)-f(b-1)).
\end{aligned}
\]
Every intermediate term \(f(a+1), f(a+2), \dots, f(b-1)\) appears once
with a positive sign and once with a negative sign, so they all cancel.
The remaining terms are \(f(b)\) and \(-f(a)\), giving \(f(b)-f(a)\).
\end{proof}

\section*{Cancellation Pattern}

The cancellation can be seen explicitly:
\begin{center}
\begin{tikzpicture}[
node distance=1.6cm,
every node/.style={font=\small}
]
\node (t1) {$f(a+1)-f(a)$};
\node (t2) [below of=t1] {$f(a+2)-f(a+1)$};
\node (t3) [below of=t2] {$f(a+3)-f(a+2)$};
\node (vdots) [below of=t3] {$\vdots$};
\node (t4) [below of=vdots] {$f(b)-f(b-1)$};

\node[right=3.2cm of t1] {$-f(a)$};
\node[right=3.2cm of t2] {$\cancel{f(a+1)}-\cancel{f(a+1)}$};
\node[right=3.2cm of t3] {$\cancel{f(a+2)}-\cancel{f(a+2)}$};
\node[right=3.2cm of vdots] {$\vdots$};
\node[right=3.2cm of t4] {$f(b)$};
\end{tikzpicture}
\end{center}

\begin{remark*}[Intuition]
Each intermediate term appears twice in the expanded sum: once with a
positive sign and once with a negative sign. These pairs cancel, leaving
only the first and last terms.
\end{remark*}

\section*{Operator Form}

Using the forward difference operator \(\Delta f(n) = f(n+1) - f(n)\),
the telescoping identity takes the compact form
\[
\sum_{k=a}^{b-1} \Delta f(k) = f(b) - f(a).
\]

\begin{remark*}[Discrete Fundamental Theorem of Calculus]
This identity is the discrete analogue of the Fundamental Theorem of
Calculus:
\[
\int_a^b f'(x)\,dx = f(b)-f(a).
\]
It reveals a deep relationship between the difference operator and
summation: summation is the inverse of differencing, just as integration
is the inverse of differentiation. This connection is made precise in the
next section through the concept of the \emph{discrete antiderivative}.
\end{remark*}
